\chapter{Tensors and the Cost of Multiplying Matrices}\label{linear:ch:tensors}

\begin{center}\itshape
Below we will give an algorithm which computes the coefficients of the product of two square matrices $A$ and $B$ of order $n$
with less than $4.7\cdot n^{\log7}$ arithmetical operations.
\upshape(Volker Strassen, \emph{Gaussian elimination is not optimal}, 1969)
\end{center}

To multiply two $2\times2$ matrices, the schoolbook rule computes eight products of entries: each of the four entries of $AB$ is
a sum of two products. For a century and a half nobody seriously asked whether eight was necessary; it seemed obvious. In 1969
Volker Strassen, then at the University of Zurich, set out (so the story goes) to prove that it was, and ended up with an
algorithm that uses seven. Seven instead of eight is a modest saving, but the algorithm uses only additions and
products of entries, never their commutativity, so it applies to $2\times2$ \emph{block} matrices whose blocks are themselves large
matrices. Applied recursively, it multiplies $n\times n$ matrices with about $n^{\log_27}\approx n^{2.807}$ operations instead of
$n^3$. The paper was two and a half pages long, its title was a provocation, and it started a race that is still on: the best
exponent known today is below $2.372$, and nobody knows whether $n^{2+\varepsilon}$ is possible for every $\varepsilon>0$.

The right language for this race is the language of tensors. A bilinear map, such as $(A,B)\mapsto AB$, is a three-way array of
numbers, and the number of multiplications it needs is the \emph{rank} of that array. Tensor rank generalises matrix rank, but it
is much stranger: it depends on the field, it is not semicontinuous, and computing it is NP-hard. This chapter builds the
tools (Kronecker products and the vectorisation trick, which make linear equations in matrices tractable), proves Strassen's
theorem, explores the strange geometry of tensor rank on $2\times2\times2$ arrays, and ends with the one case where everything
is as nice as for matrices: tensors with two indices, where the singular value decomposition becomes the \emph{Schmidt
decomposition} of quantum mechanics.

\begin{think}
Complex numbers are pairs of reals, and $(a+bi)(c+di)=(ac-bd)+(ad+bc)i$ costs four real multiplications. Can you do it with
three, if additions are free? (Gauss knew how.)
\end{think}

\section{Kronecker products and vectorisation}

\begin{definition}\label{linear:def:kronecker}
The \emph{Kronecker product} of $A\in M_{m,n}(\C)$ and $B\in M_{p,q}(\C)$ is the $mp\times nq$ block matrix
\[
A\otimes B=\begin{pmatrix}a_{11}B&\cdots&a_{1n}B\\\vdots&&\vdots\\a_{m1}B&\cdots&a_{mn}B\end{pmatrix}.
\]
For vectors this is $(u\otimes v)_{(i,k)}=u_iv_k$, the coordinates being ordered lexicographically.
\end{definition}

\begin{proposition}[rules of the Kronecker product]\label{linear:prop:kronecker}
Whenever the products are defined:
\begin{enumerate}[label=(\roman*)]
\item $(A\otimes B)(C\otimes D)=AC\otimes BD$ (the \emph{mixed product rule});
\item $(A\otimes B)^*=A^*\otimes B^*$, and $(A\otimes B)^{-1}=A^{-1}\otimes B^{-1}$ for invertible $A,B$;
\item if $A$ ($m\times m$) has eigenvalues $\lambda_1,\dots,\lambda_m$ and $B$ ($p\times p$) has eigenvalues $\mu_1,\dots,\mu_p$ (with
multiplicity), then $A\otimes B$ has eigenvalues $\lambda_i\mu_k$, and $A\otimes I_p+I_m\otimes B$ has eigenvalues $\lambda_i+\mu_k$;
\item $\operatorname{tr}(A\otimes B)=\operatorname{tr}A\operatorname{tr}B$, $\det(A\otimes B)=(\det A)^p(\det B)^m$, and
$\operatorname{rank}(A\otimes B)=\operatorname{rank}A\operatorname{rank}B$.
\end{enumerate}
\end{proposition}
\begin{proof}
(i) The $(i,j)$ block of the left side is $\sum_ka_{ik}B\,c_{kj}D=(AC)_{ij}BD$, the $(i,j)$ block of the right side. (ii) The
first identity is read off the blocks; the second follows from (i), since $(A\otimes B)(A^{-1}\otimes B^{-1})=I\otimes I=I$.

(iii) By Schur's theorem (Theorem~\ref{linear:thm:schur}) write $A=UTU^*$ and $B=VSV^*$ with $U,V$ unitary and $T,S$ upper triangular, with
the eigenvalues on their diagonals. By (i) and (ii), $A\otimes B=(U\otimes V)(T\otimes S)(U\otimes V)^*$ with $U\otimes V$ unitary,
and $T\otimes S$ is upper triangular (its block $(i,j)$ is $t_{ij}S$, zero for $i>j$, and its diagonal blocks $t_{ii}S$ are upper
triangular), with diagonal entries $t_{ii}s_{kk}=\lambda_i\mu_k$. Likewise $A\otimes I+I\otimes B$ is conjugate to the upper
triangular $T\otimes I+I\otimes S$, with diagonal $\lambda_i+\mu_k$.

(iv) Trace and determinant follow from (iii) ($\prod_{i,k}\lambda_i\mu_k=\prod_i\lambda_i^p\prod_k\mu_k^m$), or directly for the
trace. For the rank, use the singular value decompositions $A=U_1\Sigma_1V_1^*$, $B=U_2\Sigma_2V_2^*$: then $A\otimes B=(U_1\otimes U_2)
(\Sigma_1\otimes\Sigma_2)(V_1\otimes V_2)^*$, and $\Sigma_1\otimes\Sigma_2$ is a rectangular diagonal matrix whose non-zero entries are the
products $\sigma_i(A)\sigma_k(B)$, $\operatorname{rank}A\cdot\operatorname{rank}B$ of them.
\end{proof}

The proof of (iv) shows more: the singular values of $A\otimes B$ are the products of those of $A$ and $B$, so $\|A\otimes B\|_2=
\|A\|_2\|B\|_2$. Kronecker products thus carry eigenvalues and singular values along with them, and they explain the spectra of
many graphs at once.

\begin{example}[grids and cubes]\label{linear:exa:grid}
The \emph{Cartesian product} $G\mathbin\square H$ of two graphs has vertex set $V(G)\times V(H)$, with $(g,h)\sim(g',h')$ when
$g\sim g'$ and $h=h'$, or $g=g'$ and $h\sim h'$. Its adjacency matrix is $A_G\otimes I+I\otimes A_H$, so its eigenvalues are the sums
$\lambda_i(G)+\lambda_k(H)$. The $m\times n$ grid $P_m\mathbin\square P_n$ has eigenvalues $2\cos\frac{\pi i}{m+1}+2\cos\frac{\pi k}
{n+1}$ (the path $P_m$ has eigenvalues $2\cos\frac{\pi i}{m+1}$, Exercise~\ref{linear:ex:pathgraph}); the hypercube $Q_n$, the Cartesian
product of $n$ copies of $K_2$, has the eigenvalues $\pm1\pm1\dots\pm1$, that is $n-2t$ with multiplicity $\binom nt$, as in
Exercise~\ref{linear:ex:hypercube}. The same computation gives the eigenvalues of the discrete Laplacian on a rectangle, which every
numerical solver of the heat equation needs.
\end{example}

Kronecker products turn linear equations whose unknown is a matrix into ordinary linear systems. Let $\operatorname{vec}X\in\C^{mn}$
be the vector obtained by stacking the columns of $X\in M_{m,n}(\C)$, the first column on top.

\begin{proposition}[the vec trick]\label{linear:prop:vec}
For matrices of compatible sizes, $\operatorname{vec}(AXB)=(B^{\mathsf T}\otimes A)\operatorname{vec}X$.
\end{proposition}
\begin{proof}
The $j$-th column of $AXB$ is $AXBe_j=\sum_kb_{kj}AXe_k=\sum_k(b_{kj}A)x_k$, where $x_k$ is the $k$-th column of $X$. The $j$-th block
row of $B^{\mathsf T}\otimes A$ is $(b_{1j}A,\dots,b_{nj}A)$, and applied to $\operatorname{vec}X=(x_1;\dots;x_n)$ it gives the same
vector.
\end{proof}

\begin{theorem}[Sylvester's equation]\label{linear:thm:sylvester}
Let $A\in M_m(\C)$, $B\in M_n(\C)$. The equation $AX-XB=C$ has a unique solution $X\in M_{m,n}(\C)$ for every $C$ if and only if $A$
and $B$ have no common eigenvalue.
\end{theorem}
\begin{proof}
By Proposition~\ref{linear:prop:vec}, the equation reads $(I_n\otimes A-B^{\mathsf T}\otimes I_m)\operatorname{vec}X=\operatorname{vec}C$, a
square linear system. By Proposition~\ref{linear:prop:kronecker}(iii), applied to $A$ and $-B^{\mathsf T}$ (which has the eigenvalues
$-\mu_k$ of $-B$), its matrix has eigenvalues $\lambda_i-\mu_k$. It is invertible exactly when none of these is zero.
\end{proof}

Sylvester's equation is everywhere: in control theory, in the stability of differential equations (Appendix~L
studies its symmetric cousin, the Lyapunov equation), and in the block diagonalisation of matrices (Exercise~\ref{linear:ex:roth}).

\section{Tensors and their rank}

From now on, a \emph{tensor} of format $n_1\times n_2\times n_3$ over a field $K$ is a three-way array $T=(t_{ijk})$ with entries in
$K$; the space of such arrays is $K^{n_1}\otimes K^{n_2}\otimes K^{n_3}$. For vectors $u,v,w$, the \emph{rank-one tensor}
$u\otimes v\otimes w$ has entries $u_iv_jw_k$.

\begin{definition}\label{linear:def:tensorrank}
The \emph{rank} $R_K(T)$ of a tensor is the least $r$ such that $T=\sum_{s=1}^ru_s\otimes v_s\otimes w_s$ with vectors over $K$.
\end{definition}

For two indices this is the usual rank of a matrix, since $u\otimes v$ is the rank-one matrix $uv^{\mathsf T}$. For three indices, the
easiest lower bounds come from flattening. The \emph{flattening} $T_{(1)}$ is the $n_1\times n_2n_3$ matrix with rows indexed by $i$
and columns by pairs $(j,k)$. A rank-one tensor flattens to a rank-one matrix ($u\,(v\otimes w)^{\mathsf T}$), so by subadditivity of
matrix rank,
\[
R(T)\ge\max\bigl(\operatorname{rank}T_{(1)},\operatorname{rank}T_{(2)},\operatorname{rank}T_{(3)}\bigr).
\]
For matrices the two notions agree; for tensors the inequality is often strict, as we are about to see.

Fixing the last index slices a tensor into matrices: $T$ is the pair, triple, \dots\ of $n_1\times n_2$ matrices $T_k=(t_{ijk})_{i,j}$.
The case $n_1=n_2=n_3=2$ is already full of surprises, and the following lemma decodes it.

\begin{lemma}[the slice test]\label{linear:lem:slices}
Let $T\in K^2\otimes K^2\otimes K^2$ have slices $A=T_1$ and $B=T_2$, with $A$ invertible. Then $R_K(T)\le2$ if and only if $A^{-1}B$
is diagonalisable over $K$.
\end{lemma}
\begin{proof}
Suppose $T=\sum_{s=1}^2u_s\otimes v_s\otimes w_s$ (allowing zero terms). Taking slices, $A=UD_1V^{\mathsf T}$ and $B=UD_2V^{\mathsf T}$, where
$U=(u_1\;u_2)$, $V=(v_1\;v_2)$ and $D_k=\operatorname{diag}(w_{1k},w_{2k})$. As $A$ is invertible, so are $U$, $V$ and $D_1$, and
$A^{-1}B=(V^{\mathsf T})^{-1}D_1^{-1}D_2V^{\mathsf T}$ is diagonalisable over $K$.

Conversely, let $A^{-1}B=PDP^{-1}$ with $D=\operatorname{diag}(d_1,d_2)$ and $P$ invertible over $K$. Then $A=(AP)P^{-1}$ and $B=(AP)D
P^{-1}$. Writing $a_s$ for the columns of $AP$ and $b_s^{\mathsf T}$ for the rows of $P^{-1}$, $A=\sum_sa_sb_s^{\mathsf T}$ and $B=\sum_sd_sa_s
b_s^{\mathsf T}$, which is to say $T=\sum_{s=1}^2a_s\otimes b_s\otimes(1,d_s)$.
\end{proof}

\begin{example}[the rank depends on the field]\label{linear:exa:rotation}
Let $T$ have slices $A=I$ and $B=J=\bigl(\begin{smallmatrix}0&-1\\1&0\end{smallmatrix}\bigr)$, the rotation by a quarter turn. Then
$A^{-1}B=J$ has eigenvalues $\pm i$: it is diagonalisable over $\C$ but not over $\R$. So $R_\C(T)=2$ but $R_\R(T)\ge3$. In fact
$R_\R(T)=3$, as the decomposition
\[
T=(1,1)\otimes(0,1)\otimes(1,0)+(-1,1)\otimes(1,0)\otimes(0,1)+(1,0)\otimes(1,-1)\otimes(1,1)
\]
shows (check the two slices). The rank of a real tensor can drop when complex numbers are allowed, which never happens for
matrices. The same tensor appears in disguise in Exercise~\ref{linear:ex:complexmult}: it is the multiplication of complex numbers.
\end{example}

\begin{example}[rank is not semicontinuous]\label{linear:exa:wtensor}
The tensor
\[
W=e_1\otimes e_1\otimes e_2+e_1\otimes e_2\otimes e_1+e_2\otimes e_1\otimes e_1
\]
(Figure~\ref{linear:fig:wtensor}; physicists call it the W~state of three qubits) has slices $A=\bigl(\begin{smallmatrix}0&1\\1&0
\end{smallmatrix}\bigr)$ and $B=\bigl(\begin{smallmatrix}1&0\\0&0\end{smallmatrix}\bigr)$, and $A^{-1}B=\bigl(\begin{smallmatrix}0&0\\1&0
\end{smallmatrix}\bigr)$ is nilpotent and non-zero, hence not diagonalisable over any field. So $R(W)=3$. But for $\varepsilon\ne0$,
\[
W_\varepsilon=\frac1\varepsilon\Bigl((e_1+\varepsilon e_2)^{\otimes3}-e_1^{\otimes3}\Bigr)=W+\varepsilon\,(e_1\otimes e_2\otimes e_2+
e_2\otimes e_1\otimes e_2+e_2\otimes e_2\otimes e_1)+\varepsilon^2e_2^{\otimes3}
\]
has rank at most $2$, and $W_\varepsilon\to W$. A limit of tensors of rank $2$ has rank $3$.
\end{example}

\begin{figure}[ht]
\centering
\fig{linear-ch10-wtensor}
\caption{The tensor $W$ as a cube of entries. Its front slice $A$ is a permutation matrix and its back slice $B$ has a single
$1$; since $A^{-1}B$ is a non-zero nilpotent matrix, $W$ has rank $3$, yet it is a limit of tensors of rank $2$.}
\label{linear:fig:wtensor}
\end{figure}

For matrices, the set of matrices of rank at most $r$ is closed (it is defined by the vanishing of minors), and the Eckart--Young
theorem (Theorem~\ref{linear:thm:eckartyoung}) gives a best approximation of rank at most $r$. Example~\ref{linear:exa:wtensor} shows that both
facts fail for tensors: the set of tensors of rank at most $2$ is not closed, and $W$ has no best approximation of rank at most
$2$, since its distance to that set is $0$ but not attained (Exercise~\ref{linear:ex:noapprox}). Numerical methods that fit low-rank tensor
models to data meet this phenomenon as ``degenerate'' solutions, in which several rank-one terms grow without bound while
nearly cancelling. The \emph{border rank} $\underline R(T)$ is the least $r$ such that $T$ is a limit of tensors of rank $r$; so
$\underline R(W)=2<3=R(W)$.

\section{Strassen's algorithm}

The product of an $n\times m$ matrix by an $m\times p$ matrix is a bilinear map, and the right way to measure its cost is to count
multiplications. A \emph{bilinear algorithm} of length $r$ for this product consists of linear forms $\alpha_1,\dots,\alpha_r$ in the
entries of $A$, linear forms $\beta_1,\dots,\beta_r$ in the entries of $B$, and constant matrices $\Gamma_1,\dots,\Gamma_r$, such that
\[
AB=\sum_{s=1}^r\alpha_s(A)\,\beta_s(B)\,\Gamma_s\qquad\text{for all }A,B.
\]
It computes $AB$ with $r$ multiplications, all other operations being additions and multiplications by constants.

\begin{definition}\label{linear:def:mmtensor}
The \emph{matrix multiplication tensor} $\langle n,m,p\rangle$ is the tensor in $K^{nm}\otimes K^{mp}\otimes K^{pn}$ given by
\[
\langle n,m,p\rangle=\sum_{i=1}^n\sum_{j=1}^m\sum_{k=1}^pe_{ij}\otimes e_{jk}\otimes e_{ki},
\]
where $e_{ij}$ denotes the matrix unit. As a trilinear form it is $(A,B,C)\mapsto\operatorname{tr}(ABC)$.
\end{definition}

\begin{theorem}\label{linear:thm:bilinear}
The shortest bilinear algorithm for $n\times m$ by $m\times p$ matrix products has length $R(\langle n,m,p\rangle)$.
\end{theorem}
\begin{proof}
A decomposition $\langle n,m,p\rangle=\sum_{s=1}^ru_s\otimes v_s\otimes w_s$ means, as trilinear forms, that
\[
\operatorname{tr}(ABC)=\sum_s\alpha_s(A)\beta_s(B)\gamma_s(C)\qquad\text{for all }A,B,C,
\]
with $\alpha_s(A)=\sum u_{s,ij}a_{ij}$ and so on. Since $\operatorname{tr}(ABC)=\sum_{i,k}(AB)_{ik}c_{ki}$, comparing the coefficients of $c_{ki}$ gives $(AB)_{ik}=\sum_s\alpha_s(A)\beta_s(B)w_{s,ki}$: a bilinear
algorithm of length $r$ with $(\Gamma_s)_{ik}=w_{s,ki}$. The argument can be read backwards: a bilinear algorithm gives
$\operatorname{tr}(ABC)=\sum_s\alpha_s(A)\beta_s(B)\operatorname{tr}(\Gamma_sC)$, a decomposition of length $r$.
\end{proof}

The flattening bound gives $R(\langle n,n,n\rangle)\ge n^2$ (Exercise~\ref{linear:ex:flattening}); the schoolbook algorithm shows
$R(\langle n,n,n\rangle)\le n^3$. For $n=2$ the bounds are $4$ and $8$.

\begin{theorem}[Strassen, 1969]\label{linear:thm:strassen}
$R(\langle2,2,2\rangle)\le7$. Explicitly, for $2\times2$ matrices $A,B$ (whose entries may lie in any ring, commutative or not),
\[
\begin{aligned}
M_1&=(a_{11}+a_{22})(b_{11}+b_{22}),&M_5&=(a_{11}+a_{12})b_{22},\\
M_2&=(a_{21}+a_{22})b_{11},&M_6&=(a_{21}-a_{11})(b_{11}+b_{12}),\\
M_3&=a_{11}(b_{12}-b_{22}),&M_7&=(a_{12}-a_{22})(b_{21}+b_{22}),\\
M_4&=a_{22}(b_{21}-b_{11}),
\end{aligned}
\]
and then
\[
AB=\begin{pmatrix}M_1+M_4-M_5+M_7&M_3+M_5\\M_2+M_4&M_1-M_2+M_3+M_6\end{pmatrix}.
\]
\end{theorem}
\begin{proof}
Expand, keeping every $a$ to the left of every $b$. For instance $M_3+M_5=a_{11}b_{12}-a_{11}b_{22}+a_{11}b_{22}+a_{12}b_{22}=
a_{11}b_{12}+a_{12}b_{22}$, the $(1,2)$ entry of $AB$; likewise $M_2+M_4=a_{21}b_{11}+a_{22}b_{21}$. For the $(1,1)$ entry,
\[
\begin{aligned}
M_1+M_4-M_5+M_7={}&a_{11}b_{11}+a_{11}b_{22}+a_{22}b_{11}+a_{22}b_{22}+a_{22}b_{21}-a_{22}b_{11}\\
&-a_{11}b_{22}-a_{12}b_{22}+a_{12}b_{21}+a_{12}b_{22}-a_{22}b_{21}-a_{22}b_{22}\\
={}&a_{11}b_{11}+a_{12}b_{21},
\end{aligned}
\]
and the $(2,2)$ entry is checked in the same way (Exercise~\ref{linear:ex:strassen22}). No step uses $ab=ba$.
\end{proof}

How Strassen found these formulas he never said in print; they can be derived by hand from the symmetries of the tensor, and
a numerical calculation rediscovers decompositions of length $7$ numerically. What matters is the recursion.

\begin{theorem}[the exponent of matrix multiplication]\label{linear:thm:omega}
If $R(\langle k,k,k\rangle)\le r$ with $r>k^2$, then $n\times n$ matrices can be multiplied with $O(n^{\log_kr})$ arithmetic operations.
In particular $\omega\le\log_27<2.808$, where $\omega$ is the infimum of the exponents $\tau$ for which $O(n^\tau)$ operations suffice.
\end{theorem}
\begin{proof}
A bilinear algorithm of length $r$ for $k\times k$ matrices is an identity between non-commutative polynomials, since every product
in it has an entry of $A$ on the left and an entry of $B$ on the right (Theorem~\ref{linear:thm:bilinear}). So it applies to $k\times k$
block matrices whose blocks are $k^{s-1}\times k^{s-1}$ matrices: one product of size $k^s$ costs $r$ products of size $k^{s-1}$, plus
a fixed number $c$ of block additions and scalings, each costing $k^{2(s-1)}$ operations. With $T(s)$ the cost at size $k^s$,
\[
T(s)\le rT(s-1)+c\,k^{2s},\qquad\text{so}\qquad T(s)\le r^sT(0)+c\sum_{j=0}^{s-1}r^jk^{2(s-j)}=r^s\Bigl(T(0)+c\sum_{j=0}^{s-1}\Bigl(\frac{k^2}r
\Bigr)^{s-j}\Bigr).
\]
As $k^2/r<1$, the last sum is bounded by a constant, so $T(s)=O(r^s)=O(N^{\log_kr})$ with $N=k^s$. A matrix of any size $n$ is padded
with zeros to size $N=k^s$, the least power with $N\ge n$; then $N<kn$, and the cost is $O((kn)^{\log_kr})=O(n^{\log_kr})$.
\end{proof}

\begin{figure}[ht]
\centering
\fig{linear-ch10-omega}
\caption{Upper bounds for the exponent $\omega$ of matrix multiplication, from Strassen's $\log_27\approx2.807$ (1969) to $2.3713$
(2024). Since Coppersmith and Winograd (1990), all progress has come from the same ``laser method'', refined: $2.3755$, then
$2.3729$ (2012), $2.37286$ (2014), and $2.37134$ (2024). The lower bound is still $\omega\ge2$.}
\label{linear:fig:omega}
\end{figure}

\begin{remark}
Strassen's seven is optimal: $R(\langle2,2,2\rangle)=7$, proved by Shmuel Winograd in 1971 (Hopcroft and Kerr had proved it the same
year for algorithms that do not use commutativity). For $3\times3$ matrices only $19\le R(\langle3,3,3\rangle)\le23$ is known, the upper
bound being Julian Laderman's algorithm of 1976 and the lower bound Markus Bl\"aser's of 2003. In 2022 a team at DeepMind trained
a reinforcement learning agent, AlphaTensor, to play a ``game'' whose moves subtract rank-one tensors; it found a decomposition of
$\langle4,4,4\rangle$ of length $47$ over the field with two elements, beating the $49$ of Strassen's algorithm applied twice. The
asymptotic records of Figure~\ref{linear:fig:omega} use border rank and tensors other than $\langle n,n,n\rangle$ (Bini's and Sch\"onhage's
ideas), and they are galactic: the hidden constants make them useless for any matrix that fits in the universe. Strassen's
algorithm, on the other hand, is used in practice for matrices of size a few thousand.
\end{remark}

\section{Two indices: the Schmidt decomposition}

Tensors with two indices behave as well as matrices, because they \emph{are} matrices. The Hilbert space of a quantum system made of
two parts, of dimensions $m$ and $n$, is $\C^m\otimes\C^n=\C^{mn}$, with orthonormal basis $e_i\otimes e_j$; a unit vector $\psi=\sum_{i,j}
\psi_{ij}\,e_i\otimes e_j$ is a \emph{state}. It is a \emph{product state} if $\psi=a\otimes b$, and \emph{entangled} otherwise.

\begin{theorem}[Schmidt decomposition]\label{linear:thm:schmidt}
Every $\psi\in\C^m\otimes\C^n$ can be written
\[
\psi=\sum_{s=1}^r\sigma_s\,a_s\otimes b_s,
\]
with $(a_s)$ orthonormal in $\C^m$, $(b_s)$ orthonormal in $\C^n$, and $\sigma_1\ge\dots\ge\sigma_r>0$. The numbers $\sigma_s$ are the
singular values of the coefficient matrix $\Psi=(\psi_{ij})$, and $r=\operatorname{rank}\Psi$. In particular, $\psi$ is a product
state if and only if $\operatorname{rank}\Psi=1$, and $\sum_s\sigma_s^2=\|\psi\|^2$.
\end{theorem}
\begin{proof}
The map $\psi\mapsto\Psi$ is a linear isometry from $\C^m\otimes\C^n$ to $M_{m,n}(\C)$ with the Frobenius norm, and it sends $a\otimes b$
to $ab^{\mathsf T}$. Take a singular value decomposition $\Psi=\sum_s\sigma_su_sv_s^*$ (Theorem~\ref{linear:thm:svd}), and set $a_s=u_s$,
$b_s=\overline{v_s}$; the $b_s$ are orthonormal because the $v_s$ are. Then $\Psi=\sum_s\sigma_sa_sb_s^{\mathsf T}$, which is the
claimed decomposition. A product $a\otimes b\ne0$ corresponds to a matrix of rank one, and conversely. Finally $\|\psi\|^2=
\|\Psi\|_F^2=\sum_s\sigma_s^2$.
\end{proof}

The Schmidt coefficients measure entanglement. If a state $\psi$ is shared by Alice (first factor) and Bob (second factor), all
that Alice can observe is described by the \emph{reduced density matrix} $\rho_A=\Psi\Psi^*$, a positive semidefinite matrix of trace
$1$; Bob's is $\rho_B=(\Psi^*\Psi)^{\mathsf T}$. By Theorem~\ref{linear:thm:schmidt} both have the eigenvalues $\sigma_s^2$ (and zeros), so
Alice and Bob see the same amount of randomness, measured by the \emph{entanglement entropy}
\[
E(\psi)=-\sum_s\sigma_s^2\log_2\sigma_s^2.
\]
It is $0$ for product states and at most $\log_2\min(m,n)$, with equality when all the $\sigma_s$ are equal. The \emph{Bell state}
$\frac1{\sqrt2}(e_1\otimes e_1+e_2\otimes e_2)$ has $\Psi=\frac1{\sqrt2}I$ and one bit of entanglement, the maximum for two qubits.
Example~\ref{linear:exa:wtensor} shows that for three parties no such decomposition exists in general: the W~state is not a sum of two
products, let alone of orthogonal ones.

\section{A gem: three multiplications for complex numbers}

The question of the opening ``Pause and think'' has a cheerful answer. To compute $(a+bi)(c+di)$, form
\[
P_1=ac,\qquad P_2=bd,\qquad P_3=(a+b)(c+d),
\]
and return $(P_1-P_2)+(P_3-P_1-P_2)\,i$. Three real multiplications instead of four, at the price of extra additions. And three is
optimal: the bilinear map $\C\times\C\to\C$, written over $\R$, is the tensor with slices $\bigl(\begin{smallmatrix}1&0\\0&-1\end{smallmatrix}
\bigr)$ (real part) and $\bigl(\begin{smallmatrix}0&1\\1&0\end{smallmatrix}\bigr)$ (imaginary part), and the slice test shows that its
real rank is $3$ (Exercise~\ref{linear:ex:complexmult}). Over $\C$ the rank drops to $2$, because over $\C$ the problem is really two
independent multiplications in disguise. This trick, often attributed to Gauss, is how large-integer and polynomial
multiplication beat the schoolbook method: applied recursively to $(a+bX)(c+dX)$ it becomes Karatsuba's algorithm (1960), with
$O(n^{\log_23})$ digit operations instead of $O(n^2)$.

\begin{beauty}
Matrix rank is a well-tamed number: it is the same over every field containing the entries, it can only drop in a limit, and the
singular value decomposition delivers the best approximation of any rank. Add a third index, and all three facts break. Yet
that wilder notion of rank is exactly the number of multiplications that an algorithm needs, so the ancient task of multiplying
matrices is secretly a question about the geometry of tensors, and one of the great open problems of mathematics, ``is $\omega=2$?'',
asks about the rank of one explicit family of arrays of zeros and ones.
\end{beauty}

\section{Exercises}

\begin{exercise}[$\star$]\label{linear:ex:kroncommute}
Show that there is a permutation matrix $P$ (the \emph{commutation matrix}) such that $B\otimes A=P(A\otimes B)P^{\mathsf T}$ for all
$A\in M_m(\C)$, $B\in M_n(\C)$. Deduce that $A\otimes B$ and $B\otimes A$ have the same spectrum.
\end{exercise}

\begin{exercise}[$\star$]\label{linear:ex:kronnormal}
Show that $A\otimes B$ is normal (respectively Hermitian positive definite, unitary) when $A$ and $B$ are.
\end{exercise}

\begin{exercise}[$\star\star$]\label{linear:ex:roth}
Let $A\in M_m(\C)$ and $B\in M_n(\C)$ have no common eigenvalue. Show that $\bigl(\begin{smallmatrix}A&C\\0&B\end{smallmatrix}\bigr)$ is similar to
$\bigl(\begin{smallmatrix}A&0\\0&B\end{smallmatrix}\bigr)$ for every $C$.
\end{exercise}

\begin{exercise}[$\star$]\label{linear:ex:flattening}
Compute the flattening $\langle n,n,n\rangle_{(1)}$ and show that its rank is $n^2$. Deduce that $R(\langle n,n,n\rangle)\ge n^2$.
\end{exercise}

\begin{exercise}[$\star$]\label{linear:ex:strassen22}
Check the $(2,2)$ entry in Strassen's formulas, and count the additions: show that one product of size $2n$ costs $7$ products and
$18$ additions of size $n$.
\end{exercise}

\begin{exercise}[$\star\star$]\label{linear:ex:strassencount}
Let $T(n)$ be the number of arithmetic operations (additions and multiplications of numbers) of Strassen's algorithm, applied
recursively down to $1\times1$ matrices, for $n=2^s$. Show that $T(n)=7n^{\log_27}-6n^2$, and compare with the $2n^3-n^2$ operations of
the schoolbook algorithm: from which $n$ on is Strassen better?
\end{exercise}

\begin{exercise}[$\star\star$]\label{linear:ex:complexmult}
Show that the real rank of the tensor of complex multiplication (slices $\bigl(\begin{smallmatrix}1&0\\0&-1\end{smallmatrix}\bigr)$ and
$\bigl(\begin{smallmatrix}0&1\\1&0\end{smallmatrix}\bigr)$) is $3$ and its complex rank is $2$. Write the complex decomposition explicitly.
\end{exercise}

\begin{exercise}[$\star$]\label{linear:ex:noapprox}
Show that $W$ has no best approximation by tensors of rank at most $2$: for every tensor $S$ of rank at most $2$, $\|W-S\|>0$, while
$\inf_S\|W-S\|=0$ (Frobenius norms). Show also that $\|W_\varepsilon\|$ stays bounded but the two rank-one terms in the definition of
$W_\varepsilon$ have norm $1/\varepsilon$.
\end{exercise}

\begin{exercise}[$\star\star$]\label{linear:ex:waring}
Show that $24xyz=(x+y+z)^3-(-x+y+z)^3-(x-y+z)^3-(x+y-z)^3$. Interpret this as a decomposition of the symmetric tensor
$\sum_{\text{perm}}e_1\otimes e_2\otimes e_3$ into four rank-one symmetric tensors $\ell\otimes\ell\otimes\ell$.
\end{exercise}

\begin{exercise}[$\star$]\label{linear:ex:schmidtcheck}
Find the Schmidt decomposition and the entanglement entropy of $\psi=\frac13(e_1\otimes e_1+2e_1\otimes e_2+2e_2\otimes e_1)$. Is it a
product state?
\end{exercise}

\begin{exercise}[$\star\star$]\label{linear:ex:gridlap}
The \emph{discrete Laplacian} on an $m\times n$ grid with zero boundary values is $L=L_m\otimes I_n+I_m\otimes L_n$, where $L_k$ is the
$k\times k$ tridiagonal matrix with $2$ on the diagonal and $-1$ beside it. Find its eigenvalues and its condition number, and show that
the latter grows like $n^2$ for an $n\times n$ grid.
\end{exercise}

\begin{exercise}[$\star\star\star$]\label{linear:ex:border}
Show that every tensor in $\C^2\otimes\C^2\otimes\C^2$ is a limit of tensors of rank at most $2$: the border rank of every
$2\times2\times2$ complex tensor is at most $2$. (Start with the tensors whose first slice is invertible.)
\end{exercise}

\section{Hints}

\begin{hint}{linear:ex:kroncommute}
$P$ sends $e_i\otimes e_k$ to $e_k\otimes e_i$. Check the identity on rank-one matrices $A=e_ie_j^{\mathsf T}$, $B=e_ke_l^{\mathsf T}$.
\end{hint}
\begin{hint}{linear:ex:kronnormal}
Use the mixed product rule, and Proposition~\ref{linear:prop:kronecker}(iii) for positivity.
\end{hint}
\begin{hint}{linear:ex:roth}
Conjugate by $\bigl(\begin{smallmatrix}I&X\\0&I\end{smallmatrix}\bigr)$ and use Theorem~\ref{linear:thm:sylvester}.
\end{hint}
\begin{hint}{linear:ex:flattening}
The columns of the flattening are indexed by pairs $(e_{jk},e_{ki})$; which rows do they hit?
\end{hint}
\begin{hint}{linear:ex:strassen22}
$M_1-M_2+M_3+M_6$: write out the twelve terms.
\end{hint}
\begin{hint}{linear:ex:strassencount}
$T(n)=7T(n/2)+18(n/2)^2$ and $T(1)=1$.
\end{hint}
\begin{hint}{linear:ex:complexmult}
Apply Lemma~\ref{linear:lem:slices}: $A^{-1}B$ is a rotation. For the complex decomposition, the eigenvectors of the rotation are $(1,\mp i)$.
\end{hint}
\begin{hint}{linear:ex:noapprox}
Rank at most $2$ is a property that $W$ does not have, and $W_\varepsilon$ has.
\end{hint}
\begin{hint}{linear:ex:waring}
Expand. For the interpretation, identify a cubic form $\sum t_{ijk}x_ix_jx_k$ with a symmetric tensor.
\end{hint}
\begin{hint}{linear:ex:schmidtcheck}
Compute the singular values of $\frac13\bigl(\begin{smallmatrix}1&2\\2&0\end{smallmatrix}\bigr)$.
\end{hint}
\begin{hint}{linear:ex:gridlap}
The eigenvalues of $L_k$ are $2-2\cos\frac{\pi j}{k+1}$ (Exercise~\ref{linear:ex:pathgraph}).
\end{hint}
\begin{hint}{linear:ex:border}
Perturb $B$ so that $A^{-1}B$ gets two distinct eigenvalues.
\end{hint}

\section{Solutions}

\begin{solution}{linear:ex:kroncommute}
\idea Swapping the two tensor factors is a permutation of coordinates, and it conjugates $A\otimes B$ into $B\otimes A$.
\writeup Let $P$ be the permutation matrix with $P(e_i\otimes e_k)=e_k\otimes e_i$ ($i\le m$, $k\le n$). For all vectors,
$P(x\otimes y)=y\otimes x$ by linearity. Then $P(A\otimes B)(x\otimes y)=P(Ax\otimes By)=By\otimes Ax=(B\otimes A)(y\otimes x)=(B\otimes A)
P(x\otimes y)$. The vectors $x\otimes y$ span $\C^{mn}$, so $P(A\otimes B)=(B\otimes A)P$, and $P^{-1}=P^{\mathsf T}$. Similar matrices
have the same spectrum.
\end{solution}

\begin{solution}{linear:ex:kronnormal}
\idea Every property here is an identity between products and adjoints, and the mixed product rule transports such identities.
\writeup If $AA^*=A^*A$ and $BB^*=B^*B$, then $(A\otimes B)(A\otimes B)^*=AA^*\otimes BB^*=A^*A\otimes B^*B=(A\otimes B)^*(A\otimes B)$.
Unitary: $(A\otimes B)(A\otimes B)^*=AA^*\otimes BB^*=I\otimes I=I$. Hermitian positive definite: $A\otimes B$ is Hermitian, and its eigenvalues
$\lambda_i\mu_k$ are positive.
\end{solution}

\begin{solution}{linear:ex:roth}
\idea Kill the off-diagonal block with a block shear; the condition that it can be killed is a Sylvester equation.
\writeup For any $X$,
\[
\begin{pmatrix}I&X\\0&I\end{pmatrix}\begin{pmatrix}A&C\\0&B\end{pmatrix}\begin{pmatrix}I&-X\\0&I\end{pmatrix}=\begin{pmatrix}A&C-AX+XB\\0&B
\end{pmatrix}.
\]
By Theorem~\ref{linear:thm:sylvester} there is $X$ with $AX-XB=C$, and then the right side is block diagonal. (Roth proved in 1952 that
the converse also holds for each single $C$: the two matrices are similar if and only if $AX-XB=C$ has a solution, with no
condition on the spectra.)
\end{solution}

\begin{solution}{linear:ex:flattening}
\idea A flattening that turns out to be a permutation matrix has full rank.
\writeup The flattening has rows indexed by $e_{ij}$ ($n^2$ of them) and columns by pairs $(e_{j'k},e_{k'i'})$ ($n^4$ of them), with an
entry $1$ exactly when $j'=j$, $k'=k$ and $i'=i$ for some $k$. So the row of $e_{ij}$ has its ones in the columns $(e_{jk},e_{ki})$,
$k=1,\dots,n$, and different rows have disjoint supports. The $n^2$ rows are non-zero with disjoint supports, hence independent:
the rank is $n^2$. So $R(\langle n,n,n\rangle)\ge n^2$.
\end{solution}

\begin{solution}{linear:ex:strassen22}
\idea Bookkeeping, with signs.
\writeup $M_1-M_2+M_3+M_6$ equals
\[
\begin{aligned}
&(a_{11}b_{11}+a_{11}b_{22}+a_{22}b_{11}+a_{22}b_{22})-(a_{21}b_{11}+a_{22}b_{11})+(a_{11}b_{12}-a_{11}b_{22})\\
&+(a_{21}b_{11}+a_{21}b_{12}-a_{11}b_{11}-a_{11}b_{12})=a_{21}b_{12}+a_{22}b_{22},
\end{aligned}
\]
the $(2,2)$ entry of $AB$. Additions: forming the factors of $M_1,\dots,M_7$ costs $2+1+1+1+1+2+2=10$ additions or subtractions of
blocks, and assembling $AB$ costs $3+1+1+3=8$: in total $18$.
\end{solution}

\begin{solution}{linear:ex:strassencount}
\idea Solve the linear recurrence by guessing the shape $\alpha7^s+\beta4^s$.
\writeup With $n=2^s$, $T(2^s)=7T(2^{s-1})+18\cdot4^{s-1}$, $T(1)=1$. Try $T(2^s)=\alpha7^s+\beta4^s$: then $\beta4^s=7\beta4^{s-1}+
\frac{18}44^s$, so $\beta=\frac74\beta+\frac92$ and $\beta=-6$; and $T(1)=\alpha-6=1$ gives $\alpha=7$. Since $7^s=n^{\log_27}$ and $4^s=n^2$,
$T(n)=7n^{\log_27}-6n^2$. Comparing with $2n^3-n^2$: for $n=2^9=512$, $T=7\cdot7^9-6\cdot512^2\approx2.81\cdot10^8$ against $2.68\cdot10^8$,
and for $n=2^{10}$, $1.97\cdot10^9$ against $2.15\cdot10^9$. So full recursion pays only from $n=1024$ on. In practice one stops the
recursion at blocks of size about $64$ and multiplies those the usual way, which makes Strassen useful much earlier.
\end{solution}

\begin{solution}{linear:ex:complexmult}
\idea The slice test turns the question into: is a rotation diagonalisable?
\writeup Writing $(a+bi)(c+di)=(ac-bd)+(ad+bc)i$, the tensor has slices $A=\bigl(\begin{smallmatrix}1&0\\0&-1\end{smallmatrix}\bigr)$
(coefficients of $a_ic_j$ in the real part) and $B=\bigl(\begin{smallmatrix}0&1\\1&0\end{smallmatrix}\bigr)$ (imaginary part). Then
$A^{-1}B=\bigl(\begin{smallmatrix}0&1\\-1&0\end{smallmatrix}\bigr)$, with eigenvalues $\pm i$: not diagonalisable over $\R$, so $R_\R\ge3$
by Lemma~\ref{linear:lem:slices}, and Gauss's three products give $R_\R\le3$. Over $\C$, the lemma gives $R_\C\le2$; explicitly, with $z=a+bi$
and $w=c+di$, $zw$ and $\bar z\bar w$ are each one complex product, and
\[
\operatorname{Re}(zw)=\tfrac12\bigl(zw+\bar z\bar w\bigr),\qquad\operatorname{Im}(zw)=\tfrac1{2i}\bigl(zw-\bar z\bar w\bigr),
\]
where $z=a+ib$, $\bar z=a-ib$ are linear forms in $(a,b)$ and similarly for $w$. So $T=(1,i)\otimes(1,i)\otimes\frac12(1,-i)+(1,-i)\otimes
(1,-i)\otimes\frac12(1,i)$, of complex rank $2$ (and not $1$, since $A$ has rank $2$).
\end{solution}

\begin{solution}{linear:ex:noapprox}
\idea A closed set always contains its limit points; the set of rank-$\le2$ tensors is not closed, and $W$ is a missing limit point.
\writeup If $\|W-S\|=0$ then $S=W$, which has rank $3$; so $\|W-S\|>0$ for every $S$ of rank at most $2$. But $W_\varepsilon$ has rank at
most $2$ and $\|W-W_\varepsilon\|\to0$, so the infimum is $0$, not attained. The norm of $W_\varepsilon$ tends to $\|W\|=\sqrt3$, while
$\|\frac1\varepsilon(e_1+\varepsilon e_2)^{\otimes3}\|=\frac{(1+\varepsilon^2)^{3/2}}{|\varepsilon|}$ and $\|\frac1\varepsilon e_1^{\otimes3}\|=\frac1
{|\varepsilon|}$ blow up: the two terms nearly cancel.
\end{solution}

\begin{solution}{linear:ex:waring}
\idea Expand by the multinomial formula; the cubes of the variables and the terms $x^2y$ cancel by symmetry, only $xyz$ survives.
\writeup In $(\pm x\pm y\pm z)^3$ the coefficient of $xyz$ is $6$ times the product of the three signs, the coefficient of $x^3$ is the
cube of the sign of $x$, and that of $x^2y$ is $3$ times the sign of $y$. Over the four sign patterns $(+,+,+)$ with weight $+1$ and
$(-,+,+),(+,-,+),(+,+,-)$ with weight $-1$: the $x^3$ terms give $1-(-1)-1-1=0$, the $x^2y$ terms give $3(1-1-(-1)-1)=0$, and the
$xyz$ terms give $6(1+1+1+1)=24$. A cubic form $f(x)=\sum t_{ijk}x_ix_jx_k$ with $t$ symmetric is determined by $t$; the form $\ell(x)^3$
corresponds to $\ell\otimes\ell\otimes\ell$, and $6x_1x_2x_3$ to the sum of the six tensors $e_{\pi(1)}\otimes e_{\pi(2)}\otimes e_{\pi(3)}$. So
that sum equals $\frac14\bigl(\ell_0^{\otimes3}-\ell_1^{\otimes3}-\ell_2^{\otimes3}-\ell_3^{\otimes3}\bigr)$ with $\ell_0=(1,1,1)$, $\ell_1=(-1,1,1)$,
$\ell_2=(1,-1,1)$, $\ell_3=(1,1,-1)$.
\end{solution}

\begin{solution}{linear:ex:schmidtcheck}
\idea Schmidt coefficients are singular values; for a symmetric matrix they are the absolute values of the eigenvalues.
\writeup $\Psi=\frac13\bigl(\begin{smallmatrix}1&2\\2&0\end{smallmatrix}\bigr)$ is symmetric with eigenvalues $\frac13\cdot\frac{1\pm\sqrt{17}}2$,
so $\sigma_1=\frac{1+\sqrt{17}}6\approx0.854$ and $\sigma_2=\frac{\sqrt{17}-1}6\approx0.520$, and
$\sigma_1^2+\sigma_2^2=\frac{(1+\sqrt{17})^2+(\sqrt{17}-1)^2}{36}=1$, as it must since $\|\psi\|^2=\frac{1+4+4}9=1$. The vectors $a_s$ are the eigenvectors $u_s$ of $\Psi$, and $b_s=\pm u_s$
(with the sign of the eigenvalue). The entropy is $-\sum\sigma_s^2\log_2\sigma_s^2\approx0.843$ bits. The state is entangled, since
$\operatorname{rank}\Psi=2$.
\end{solution}

\begin{solution}{linear:ex:gridlap}
\idea Kronecker sums add eigenvalues; the extreme eigenvalues are at the two ends of the cosine.
\writeup By Proposition~\ref{linear:prop:kronecker}(iii), the eigenvalues of $L$ are
\[
\lambda_{jk}=4-2\cos\frac{\pi j}{m+1}-2\cos\frac{\pi k}{n+1},\qquad1\le j\le m,\ 1\le k\le n,
\]
all positive. $L$ is symmetric positive definite, so its condition number is $\lambda_{\max}/\lambda_{\min}$. For $m=n$,
$\lambda_{\min}=4-4\cos\frac\pi{n+1}=8\sin^2\frac\pi{2(n+1)}\sim\frac{2\pi^2}{n^2}$ and $\lambda_{\max}=4+4\cos\frac\pi{n+1}\to8$, so
$\kappa(L)\sim\frac{4n^2}{\pi^2}$. Refining the grid makes the problem quadratically worse conditioned, which is why fast Poisson solvers
use the eigenbasis (a sine transform, computed by FFT) instead of plain iterations; Appendix~L returns to this.
\end{solution}

\begin{solution}{linear:ex:border}
\idea Diagonalisable matrices are dense, so by the slice test tensors of rank at most $2$ are dense among those with an
invertible slice; and those are dense among all tensors.
\writeup \emph{Step 1: $A$ invertible.} Let $M=A^{-1}B$, and triangularise it, $M=P\bigl(\begin{smallmatrix}\lambda&x\\0&\mu\end{smallmatrix}
\bigr)P^{-1}$. Let $M_\varepsilon=P\bigl(\begin{smallmatrix}\lambda&x\\0&\mu+\varepsilon\end{smallmatrix}\bigr)P^{-1}$: for all small
$\varepsilon\ne0$ its eigenvalues $\lambda$ and $\mu+\varepsilon$ are distinct, so it is diagonalisable over $\C$. By
Lemma~\ref{linear:lem:slices} the tensor $T_\varepsilon$ with slices $A$ and $AM_\varepsilon$ has rank at most $2$, and $T_\varepsilon\to T$
because $AM_\varepsilon\to AM=B$.

\emph{Step 2: any $T$.} The polynomial $\det(A+\varepsilon I)$ in $\varepsilon$ has leading term $\varepsilon^2$, so it vanishes for at
most two values: for all small $\varepsilon\ne0$, the tensor $T'_\varepsilon$ with slices $A+\varepsilon I$ and $B$ has an invertible
first slice, and $T'_\varepsilon\to T$. Given $m\ge1$, choose $\varepsilon$ with $\|T-T'_\varepsilon\|<\frac1m$, then by Step~1 a tensor
$S_m$ of rank at most $2$ with $\|T'_\varepsilon-S_m\|<\frac1m$. Then $S_m\to T$.
\end{solution}

\section*{Sources and further reading}

\begin{itemize}
\item V.~Strassen, \emph{Gaussian elimination is not optimal}, Numerische Mathematik 13 (1969), 354--356; S.~Winograd, \emph{On
multiplication of $2\times2$ matrices}, Linear Algebra and its Applications 4 (1971), 381--388; J.~E.~Hopcroft and L.~R.~Kerr, \emph{On
minimizing the number of multiplications necessary for matrix multiplication}, SIAM Journal on Applied Mathematics 20 (1971), 30--36.
\item J.~D.~Laderman, \emph{A noncommutative algorithm for multiplying $3\times3$ matrices using 23 multiplications}, Bulletin of the
AMS 82 (1976), 126--128; M.~Bl\"aser, \emph{On the complexity of the multiplication of matrices of small formats}, Journal of
Complexity 19 (2003), 43--60.
\item D.~Coppersmith and S.~Winograd, \emph{Matrix multiplication via arithmetic progressions}, Journal of Symbolic Computation 9
(1990), 251--280; J.~Alman, R.~Duan, V.~Vassilevska Williams, Y.~Xu, Z.~Xu and R.~Zhou, \emph{More asymmetry yields faster matrix
multiplication}, 2024 (preprint, arXiv:2404.16349).
\item A.~Fawzi et al., \emph{Discovering faster matrix multiplication algorithms with reinforcement learning}, Nature 610 (2022),
47--53 (AlphaTensor).
\item P.~B\"urgisser, M.~Clausen and M.~A.~Shokrollahi, \emph{Algebraic Complexity Theory}, Springer, 1997; J.~M.~Landsberg,
\emph{Tensors: Geometry and Applications}, AMS, 2012. The standard references on tensor rank, border rank and $\omega$.
\item V.~de Silva and L.-H.~Lim, \emph{Tensor rank and the ill-posedness of the best low-rank approximation problem}, SIAM Journal
on Matrix Analysis and Applications 30 (2008), 1084--1127; C.~J.~Hillar and L.-H.~Lim, \emph{Most tensor problems are NP-hard},
Journal of the ACM 60 (2013), article 45.
\item W.~E.~Roth, \emph{The equations $AX-YB=C$ and $AX-XB=C$ in matrices}, Proceedings of the AMS 3 (1952), 392--396; R.~A.~Horn and
C.~R.~Johnson, \emph{Topics in Matrix Analysis}, Cambridge University Press, 1991, Chapter~4 (Kronecker products).
\item M.~A.~Nielsen and I.~L.~Chuang, \emph{Quantum Computation and Quantum Information}, Cambridge University Press, 2000,
Section~2.5 (Schmidt decomposition); W.~D\"ur, G.~Vidal and J.~I.~Cirac, \emph{Three qubits can be entangled in two inequivalent ways},
Physical Review A 62 (2000), 062314 (the W state).
\item A.~Karatsuba and Yu.~Ofman, \emph{Multiplication of multidigit numbers on automata}, Soviet Physics Doklady 7 (1963), 595--596.
\end{itemize}
